\documentclass[12pt]{article}

\usepackage[margin=1.25in]{geometry}
\usepackage{setspace}
\usepackage{hyperref}
\usepackage{booktabs}
\usepackage{footnote}
% microtype removed for compatibility

\hypersetup{
    colorlinks=true,
    linkcolor=blue,
    citecolor=blue,
    urlcolor=blue
}

\doublespacing

\title{The Union's 5-20 Bonds: Congressional Debates, the Gold/Greenback
Ambiguity, and Its Logic}
\author{Notes prepared with reference to Hall and Sargent (2014)}
\date{\today}

\begin{document}

\maketitle
\thispagestyle{empty}

\newpage

\tableofcontents

\newpage

%-----------------------------------------------------------------------
\section{A Note on Sources}
%-----------------------------------------------------------------------

First, a bibliographic correction: the Civil War-era debates would be
found in the \textit{Congressional Globe}, not the Annals of Congress
(which covers only 1789--1824). The Globe covered the 23rd through 42nd
Congresses (1833--1873) and is available through the Library of Congress
American Memory collection at \texttt{memory.loc.gov}. The relevant
debates are in the \textbf{37th Congress, 2nd Session} (January--July
1862). The Globe is unfortunately not easily full-text searchable through
a single interface, but the relevant volumes are accessible page by page.

%-----------------------------------------------------------------------
\section{The Legislative Package: Legal Tender Act of February 25, 1862}
%-----------------------------------------------------------------------

The 5-20 bonds were not a standalone bill; they were embedded in the
\textbf{Legal Tender Act} (formally ``An Act to authorize the Issue of
United States Notes, and for the Redemption or Funding thereof, and for
Funding the Floating Debt of the United States,'' Chapter XXXIII of the
37th Congress, 2nd Session). The actual statutory text reveals the
precise structure of the currency language:

\textbf{Section~1} authorized \$150~million in U.S.\ Notes (greenbacks),
declared legal tender for all public and private debts ``except duties on
imports and interest as aforesaid.'' It then provided that holders of
greenbacks in sums of \$50 or more could convert them into bonds ``bearing
interest at the rate of six per centum per annum, payable semi-annually,
and redeemable at the pleasure of the United States after five years, and
payable twenty years from the date thereof.'' This is the 5-20 structure.

\textbf{Section~2} authorized a broader \$500~million in coupon or
registered bonds (the main 5-20 issue), ``redeemable at the pleasure of
the United States after five years, and payable twenty years from date,
and bearing interest at the rate of six per centum per annum, payable
semi-annually.'' Notably, it says nothing about coin or greenbacks for
\textit{principal}.

\textbf{Section~5}---the critical provision---required that import duties
be paid in coin, and earmarked that coin: \textit{first} to ``the payment
in coin of the interest on the bonds and notes of the United States,'' and
\textit{second} to a sinking fund of one percent of the total debt per
year for gradual principal retirement.

The Act is therefore entirely explicit that \textbf{interest is payable in
coin}, but entirely silent about \textbf{principal}. This is the source
of the great controversy that erupted after the war.

%-----------------------------------------------------------------------
\section{The Congressional Debates}
%-----------------------------------------------------------------------

The House debate was driven by the House Ways and Means Committee under
\textbf{Thaddeus Stevens} (PA) and the bill's principal architect,
\textbf{Elbridge Spaulding} (NY). Stevens opened with the famous
declaration that ``This bill is a matter of necessity, not of choice,''
framing both the greenbacks and the bond conversion as emergency war
measures. Spaulding similarly called it ``a war measure, a necessary means
of carrying into execution the power granted in the Constitution `to raise
and support armies.'\,''

Opposition came from multiple directions. \textbf{Roscoe Conkling}~(NY)
and \textbf{Justin Morrill}~(VT) objected to the legal tender provision on
constitutional and commercial grounds. The most passionate dissent came
from \textbf{George Pendleton}~(OH) and \textbf{Clement
Vallandigham}~(OH), who prophetically warned that paper money would
inflate prices and destroy the savings of ordinary workers. Pendleton
declared:
\begin{quote}
The wit of man has never discovered a means by which paper currency can
be kept at par value, except by its speedy, cheap, certain convertibility
into gold and silver\ldots\ prices will be inflated\ldots\ incomes will
depreciate; the savings of the poor will vanish; the hoardings of the
widow will melt away; bonds, mortgages, and notes---everything of fixed
value---will lose their value.
\end{quote}
This was the same Pendleton who would later sponsor the ``Ohio Idea'' to
pay the principal of these very bonds in greenbacks.

In the \textbf{Senate Finance Committee}, the bill was managed by
\textbf{William Pitt Fessenden}~(ME), who famously expressed personal
revulsion while voting for it: it was ``of doubtful
constitutionality\ldots\ It is bad faith\ldots\ It shocks all my notions
of political, moral, and national honor\ldots\ however, to leave the
government without resources in such a crisis is not to be thought of.''
\textbf{John Sherman}~(OH) was another key Senate backer.

Chase himself initially opposed the legal tender provision and wrote to
Spaulding on January~29, 1862 warning that forcing paper on creditors
would ``unsettle credit'' and invite legal challenge---he preferred notes
redeemable in specie or backed by bonds. He then reversed course under
the pressure of war finance, writing subsequently in favor of the bill.

On the specific question of \textit{what currency the principal would be
paid in}, there appears to have been remarkably little explicit floor
debate. The silence in the statute seems to have been partly a product of
the speed and pressure under which the legislation was drafted. The coin
sinking fund of Section~5 was the device that allowed supporters to imply
gold repayment without committing to it in statute.

%-----------------------------------------------------------------------
\section{Chase's Extra-Statutory Promise}
%-----------------------------------------------------------------------

Although the statute was silent on principal repayment, Chase as Treasury
Secretary went further in his bond marketing. The key evidence comes from
Jay Cooke's later correspondence. Cooke's firm wrote that the 5-20 bonds
were ``sold to [the public] at prices which could not possibly have been
obtained but for the distinct understanding that they were payable
principal and interest in coin.'' Cooke's letter then argued that the
sinking fund clause of Section~5 itself implied this: since import duties
(paid in coin) were dedicated first to interest and second to a sinking
fund for the principal, the legislative intent was clearly gold repayment
throughout.

This was Chase's own view. Though the statute was silent, his Treasury
circulars and bond prospectuses advertised the 5-20s as redeemable in
coin, and the bonds were sold---by Cooke's army of some 2,500 agents, who
ultimately raised over \$511~million in 5-20s from patriotic small
investors across the Union---on the implicit understanding that gold
repayment applied to both coupon and principal.

The arrangement thus had a hybrid character: \textbf{statutory commitment}
for interest coupons (via Section~5), and \textbf{reputational commitment}
for principal (via Treasury marketing and Chase's personal assurances).

%-----------------------------------------------------------------------
\section{The Post-War Controversy: The Ohio Idea and the 1868 Election}
%-----------------------------------------------------------------------

The ambiguity was a time bomb. By 1867--68, with the 5-20s representing
some 70\% of the outstanding interest-bearing debt and gold trading at a
40\% premium over greenbacks, the stakes were enormous. Pendleton and the
\textit{Cincinnati Enquirer} launched the \textbf{Ohio Idea} (also known
as the Pendleton Plan): since the authorizing act only explicitly specified
coin for interest, the plain statutory language permitted---indeed
required---payment of the principal in ``lawful money,'' i.e., greenbacks.
The Democratic Party's 1868 platform (the third and fourth planks) embodied
this position, arguing that bonds ``where the law under which they were
issued does not provide that they shall be paid in coin, they ought, in
right and in justice, to be paid in the lawful money of the United States.''

The argument had genuine legal force. Under the February~25, 1862 Act,
interest alone was specified in coin; the inference from this specification
is that Congress would have said the same for principal if it had intended
gold repayment. The March~1864 act authorizing a further batch of 5-20s
and ten-forties actually \textit{did} specify both principal and interest
in coin, precisely to eliminate ambiguity about the 1862 act---which
Democrats read as retroactively confirming that the original 1862 silence
had been deliberate and meaningful.

Grant and the Republicans took the opposite position: every prior matured
U.S.\ bond had been paid in coin by practice and honor, and repudiating
this tradition would destroy the country's credit. The voters sided with
Grant in November 1868, and the \textbf{Public Credit Act of March~18,
1869}---the first legislation of the 41st Congress---resolved the
ambiguity by statute, pledging ``the faith of the United States\ldots\
solemnly\ldots\ to the payment in coin'' of all obligations not expressly
designated otherwise.

%-----------------------------------------------------------------------
\section{Were There Good Arguments for Creating the Ambiguity?}
%-----------------------------------------------------------------------

This is the most interesting analytical question. There were several
arguments that the ambiguity served a purpose---and in some cases, that
it was unavoidable.

\subsection{Political Arithmetic}

The Legal Tender Act barely passed. The House Ways and Means Committee
was initially deadlocked 4--4 on the legal tender provision itself before
a single vote changed (as Spaulding recorded in his subsequent history of
the legislation). Adding an explicit gold repayment clause for principal
would have stiffened opposition from soft-money members, particularly
Western and Midwestern Democrats whose constituencies were net debtors.
The strategic silence bought the votes needed to pass the package. In
short, the ambiguity was the price of the coalition.

\subsection{Flexibility Under Uncertainty}

During a war with uncertain outcome, locking the government into gold
repayment of \$500~million in bonds created serious fiscal risk. If Union
armies had continued to lose, the gold premium would have widened further,
and gold repayment would have been crushing. The silence preserved the
option to pay in greenbacks if necessity demanded---precisely the Lucas
and Stokey (1983) state-contingent logic that early American statesmen
generally rejected \textit{in principle} (as Hall and Sargent (2014)
document at length) but practiced inadvertently in the structure of this
legislation.

\subsection{Chase's Reputational Substitute}

Chase used his personal and institutional reputation to supply a
commitment that the statute withheld. The ambiguity in law was bridged
by administrative promise, allowing maximum legislative flexibility while
preserving investor confidence. This replicates a logic familiar from the
literature on reputation and monetary commitment: the government could not
formally commit through statute, so it committed reputationally through
Treasury marketing, with Chase's personal credibility serving as the
bonding mechanism. The fragility of this arrangement was exposed only
after the war, when Chase as Chief Justice voted to invalidate the Legal
Tender Act, implicitly renouncing the very commitment he had made as
Secretary.

\subsection{Heterogeneous Beliefs and Market Segmentation}

From an MM perspective, whether principal was paid in coin or greenbacks
of equal real value would be irrelevant. But with heterogeneous beliefs
about the future gold/greenback exchange rate---which Harrison and Kreps
(1978) showed can generate interesting asset price dynamics---the
ambiguity itself had market value. Optimists about Union prospects (who
expected greenbacks to eventually trade at par with gold) were happy to
hold the bonds without an explicit coin clause; pessimists would not have
bought the bonds at reasonable prices regardless. The ambiguity thus
efficiently segmented the market toward buyers who were relatively more
optimistic about Union prospects, functioning as a signaling and
screening device. Lowenstein (\textit{Ways and Means}, 2022) captures
Chase's intuition here, calling the interest-in-coin provision ``an
alchemy machine, a device for turning paper into specie''---the gold
coupon stream made the bonds attractive even to investors who accepted
payment in depreciating greenbacks.

\subsection{Postponing an Unresolvable Disagreement}

Hall and Sargent (2014, fn.~33) capture this logic precisely: ``Government
policies are ambiguous when decision makers cannot agree what to do now
and decide to postpone decisions until a later time when they hope that
they can agree.'' The 37th Congress could not reach consensus on gold
versus greenback principal repayment, so they passed the ambiguity
forward---hoping that subsequent events (the war's outcome, the course of
inflation, the political balance of power) would make the decision easier.
In the event, Union victory and Grant's election resolved it in favor of
gold. Had the war gone otherwise, the greenback option would have been
available as an emergency measure.

%-----------------------------------------------------------------------
\section{What Lowenstein Adds}
%-----------------------------------------------------------------------

Lowenstein's \textit{Ways and Means} (2022) provides useful narrative
texture on the Congressional politics and Chase's role. Beyond the
``alchemy machine'' metaphor for the coin-interest provision, Lowenstein
emphasizes the inegalitarian character of the overall structure: soldiers
were paid in depreciating greenbacks while bondholders received gold
coupons, a class distinction that Thaddeus Stevens---despite voting for
the bill---reportedly denounced as ``a cunning scheme.'' This class
tension is precisely what the Pendleton Plan later exploited politically
in the 1867--68 cycle: if ordinary debtors and farmers had to pay taxes
in coin to service bonds that common soldiers had purchased in greenbacks,
then paying the principal in greenbacks was not repudiation but equity.

Lowenstein also emphasizes Chase's central role in the bond marketing
operation through Jay Cooke, and how deeply Chase's personal credibility
became intertwined with the implicit gold-repayment commitment. The
irony of Chase later voting as Chief Justice to strike down the legal
tender law he had championed as Treasury Secretary is one of the
episode's most striking features.

%-----------------------------------------------------------------------
\section{Summary and Analytical Interpretation}
%-----------------------------------------------------------------------

The 5-20s are a remarkable case study in deliberate legislative
ambiguity---and in the consequences of substituting reputational for
statutory commitment. The Act of February~25, 1862 explicitly protected
bondholders' coupon income via the coin-interest clause and the dedicated
customs revenue fund, but left the principal in limbo. This was likely a
compromise between hard-money Republicans who wanted full gold commitment
and members who needed fiscal flexibility under wartime uncertainty.
Chase then informally promised what Congress had withheld, creating a
hybrid commitment structure: statutory for interest, reputational for
principal.

The post-war fight over whether the informal promise was legally binding
became one of the central monetary-political conflicts of Reconstruction,
settled only when Grant's victory in 1868 made the soft-money position
politically untenable and the Public Credit Act of 1869 converted Chase's
promise into law.

From the perspective of the analytical framework in Hall and Sargent
(2014), the episode illustrates the Bryant-Wallace (1984) price
discrimination logic operating through institutional ambiguity: different
classes of creditors received different effective commitments---soldiers
received greenbacks; domestic bondholder coupons were secured in coin by
statute; bondholder principals were secured in coin by reputation only;
foreign holders received gold throughout by implicit convention. The
ambiguity also inadvertently created something close to the Lucas and
Stokey (1983) state-contingent debt structure that early American
statesmen ideologically rejected: the principal was effectively
state-contingent on the war's outcome, mediated through the gold premium.
Fudenberg and Kreps (1987), cited in the paper for their analysis of
sustaining different reputations with different parties, provide an
analytical lens for understanding how Chase managed to maintain divergent
implicit commitments to domestic small investors (gold, via prospectus)
and to Congress (greenbacks permissible, via statutory silence).

Finally, the episode prefigures the Rogoff (1985) logic of delegating
monetary commitment to a hawk. Lincoln appointed Chase precisely because
his hard-money, Jackson-Democrat credentials would provide credibility
for an otherwise suspect financing program. The ambiguity on principal
was therefore doubly hedged: by the coin-interest clause (statutory) and
by Chase's personal reputation (extra-statutory). When Chase's
credibility was later tested---first by his Supreme Court vote, then by
the political fight of 1868---it was Grant who ultimately had to convert
the reputational commitment into the durable statutory form that the
original legislation had failed to provide.

%-----------------------------------------------------------------------
\section*{References}
%-----------------------------------------------------------------------

\begin{description}
\setlength\itemsep{0.5em}

\item Bryant, John and Neil Wallace. 1984. ``A Price Discrimination
Analysis of Monetary Policy.'' \textit{Review of Economic Studies} 51(2):
279--288.

\item Fudenberg, Drew and David M. Kreps. 1987. ``Reputation in the
Simultaneous Play of Multiple Opponents.'' \textit{Review of Economic
Studies} 54(4): 541--568.

\item Hall, George J. and Thomas J. Sargent. 2014. ``Fiscal
Discriminations in Three Wars.'' \textit{Journal of Monetary Economics}
61: 148--166.

\item Harrison, J. Michael and David M. Kreps. 1978. ``Speculative
Investor Behavior in a Stock Market with Heterogeneous Expectations.''
\textit{Quarterly Journal of Economics} 92(2): 323--336.

\item Lucas, Robert E., Jr. and Nancy L. Stokey. 1983. ``Optimal Fiscal
and Monetary Policy in an Economy Without Capital.'' \textit{Journal of
Monetary Economics} 12(1): 55--93.

\item Lowenstein, Roger. 2022. \textit{Ways and Means: Lincoln and His
Cabinet and the Financing of the Civil War}. New York: Penguin Press.

\item Mitchell, Wesley C. 1903. \textit{A History of the Greenbacks}.
Chicago: University of Chicago Press.

\item Oberholtzer, Ellis Paxson. 1907. \textit{Jay Cooke: Financier of
the Civil War}, 2 vols. Philadelphia: George W. Jacobs.

\item Rogoff, Kenneth. 1985. ``The Optimal Degree of Commitment to an
Intermediate Monetary Target.'' \textit{Quarterly Journal of Economics}
100(4): 1169--1189.

\item Spaulding, Elbridge Gerry. 1869. \textit{History of the Legal
Tender Paper Money Issued During the Great Rebellion}. Buffalo: Express
Printing Company.

\item Unger, Irwin. 1964. \textit{The Greenback Era: A Social and
Political History of American Finance, 1865--1879}. Princeton: Princeton
University Press.

\item United States. Congress. \textit{Congressional Globe}, 37th
Congress, 2nd Session (1862). Washington, D.C.: Library of Congress
American Memory. \texttt{memory.loc.gov/ammem/amlaw/lwcg.html}.

\item United States. Statutes at Large, Vol.~XII, Chap.~XXXIII (February
25, 1862). Legal Tender Act. St.\ Louis Fed FRASER:
\texttt{fraser.stlouisfed.org/title/legal-tender-act-1107}.

\end{description}

\end{document}
